% Generated draft: requires mathematical and source review.
\documentclass[11pt]{article}
\usepackage[margin=25mm]{geometry}
\usepackage{fontspec}
\setmainfont{Latin Modern Roman}
\usepackage{amsmath,amssymb,mathtools,unicode-math}
\setmathfont{Latin Modern Math}
\usepackage{enumitem,needspace,xurl,hyperref}
\setlength{\parindent}{0pt}
\setlength{\parskip}{5pt}
\setlength{\emergencystretch}{3em}
\newcommand{\N}{\mathbb{N}}
\newcommand{\Z}{\mathbb{Z}}
\newcommand{\Q}{\mathbb{Q}}
\newcommand{\R}{\mathbb{R}}
\newcommand{\C}{\mathbb{C}}
\newcommand{\F}{\mathbb{F}}
\newcommand{\A}{\mathbb{A}}
\newcommand{\PP}{\mathbb{P}}
\newcommand{\E}{\mathbb{E}}
\newcommand{\eps}{\varepsilon}
\newcommand{\dd}{\,\mathrm{d}}
\begin{document}
\section*{UnsolvedMath ID: 30006131}
\textbf{Source problem code:} catalogue entry
\subsection{Definitions and mathematical statement}
Let \(n\geq 3\) and \(m\geq 2\) be integers. Interpret a convex body \(K\subset\R^n\) as a compact convex set with nonempty interior. Write \(|\cdot|_d\) for \(d\)-dimensional Lebesgue measure and \(B_2^n=\{x\in\R^n:\|x\|_2\leq 1\}\).

The supplied original excerpt has damaged intersection notation and does not explicitly define \(D^m(K)\). For this draft, interpret the higher-order difference body using a common intersection of translates:
\[
D^m(K):=\left\{(x_1,\ldots,x_m)\in(\R^n)^m:
K\cap\bigcap_{i=1}^m(K+x_i)\neq\varnothing\right\}.
\]
Equivalently, under this interpretation,
\[
D^m(K)=\left\{(y_0-y_1,\ldots,y_0-y_m):y_0,y_1,\ldots,y_m\in K\right\}.
\]
This reconstruction of the notation requires review against the cited report. Define
\[
\delta_m(K):=\frac{|D^m(K)|_{nm}}{|K|_n^m}.
\]
An ellipsoid means a set \(a+T(B_2^n)\), where \(a\in\R^n\) and \(T:\R^n\to\R^n\) is invertible and linear.

\textbf{Conjecture attributed in the supplied record to Schneider.} For every pair of integers \(n\geq 3\), \(m\geq 2\), and every convex body \(K\subset\R^n\),
\[
\delta_m(K)\geq\delta_m(B_2^n),
\]
with equality if and only if \(K\) is an ellipsoid. Thus the minimizers are precisely the ellipsoids.

The convention that convex bodies have nonempty interior ensures that the denominator is positive; the supplied statement does not explicitly specify this convention. The supplied record labels the problem partially solved and labels its reconstructed statement as verified against the official report. These are supplied curation claims, not independent verification in this draft.
\subsection{Short English statement}
Schneider's conjecture states that, for every integer dimension \(n\geq3\) and order \(m\geq2\), ellipsoids are exactly the convex bodies minimizing the volume of the higher-order difference body divided by the \(m\)th power of the original volume. The definition of the higher-order difference body is reconstructed here from damaged source notation and requires human review.
\subsection{Sources}
\begin{itemize}
\item Local UnsolvedMath registry, permanent ID \(30006131\), title: \emph{Schneider Higher-Order Difference-Body Conjecture}; exact source problem code: \texttt{catalogue entry}. The supplied record attributes the conjecture to Schneider but supplies no full name or individual proposer. It records proposed year 2024 and states that its reconstructed statement was checked against the official Oberwolfach report; that verification claim is supplied metadata.
\item Supplied report citation: \emph{Convex Geometry and its Applications (2025)}, Oberwolfach Reports, pp.\ 3301--3376, DOI \texttt{10.4171/owr/2024/58}. \url{https://doi.org/10.4171/owr/2024/58}. The supplied background describes Oberwolfach Reports as documenting mathematical workshops and their open-problem discussions.
\item Supplied literature source: \emph{On the polar of Schneider's difference body}. \url{https://arxiv.org/abs/2503.06191}.
\item Supplied literature source: \emph{OWR 2024/58}. \url{https://ems.press/content/serial-article-files/50771}.
\item Literature assessment dated exactly \texttt{2026-08-22}: the supplied curation labels the problem partially solved and states that the conjecture remains open in general, with special cases, asymptotic results, and a related polar dual problem proved. This is supplied dated curation derived from the OWR record, not independently verified here; time-sensitive status and cited claims require independent verification.
\end{itemize}
\end{document}
